% =============================================================================
% cornell-notes-study-sheet.tex
% Style: cornell-notes
% Build: pdflatex via latexmk; minted-only.
% =============================================================================
\documentclass[10pt,letterpaper]{article}
\usepackage{cornell-notes}

\newcommand{\CornellDocumentTitle}{Discrete Mathematics Cornell Notes: Chapter 4 -- Graph Theory Basics}
\title{\CornellDocumentTitle}
\author{Documentation Team}
\date{\today}

\setDocTitle{\CornellDocumentTitle}
\setDocSubtitle{Chapter 4: Graph Theory Basics}
\setDocVersion{v1.0}
\setDocStatus{Draft}
\setDocOwner{Math Study Group}
\setDocAuthor{Documentation Team}
\setDocDate{\today}
\setDocSummary{A concise, reusable Cornell-notes example showing metadata, cue/note rows, source cues, and review callouts.}

\setCornellCollection{Discrete Mathematics Study Notes}
\setCornellUnitType{Chapter}
\setCornellUnitNumber{4}
\setCornellUnitTitle{Graph Theory Basics}
\setCornellSourceLabel{Textbook}
\setCornellSourceTitle{Discrete Mathematics and Its Applications}
\setCornellSourceAuthor{Kenneth H. Rosen}
\setCornellSourceEdition{8th Edition}
\setCornellSourceRange{pp. 645--689}
\setCornellCompanionLabel{Rapid review companion}

\begin{document}
\maketitle

\section*{Unit Overview}
Graphs model relationships among entities and support reasoning about paths,
connectivity, cycles, and optimization.

\begin{summarybox}{Review Objectives}
Explain graph terminology, compare path and cycle concepts, and apply basic
adjacency and degree properties to simple problems.
\end{summarybox}

\section{Core Concepts}
\source{pp. 645--654}
\begin{cornell}
\noterow{What is a graph?}{A graph is an ordered pair $G=(V,E)$ where $V$ is a set of vertices and $E$ is a set of edges connecting pairs of vertices.}
\noterow{Path vs. cycle}{A path is a sequence of edges connecting distinct vertices; a cycle is a closed path that starts and ends at the same vertex.}
\noterow{Degree of a vertex}{The degree of a vertex is the number of incident edges. In undirected graphs, $\sum_{v\in V}\deg(v)=2|E|$.}
\end{cornell}

\begin{exambox}{High-Yield Cue}
If exactly two vertices in a connected undirected graph have odd degree, an
Euler path exists; if none have odd degree, an Euler circuit exists.
\end{exambox}

\section{Final Checklist}
\begin{CornellChecklist}
\item I can define vertices, edges, paths, and cycles.
\item I can compute vertex degree and apply the handshake lemma.
\item I can identify whether a small graph can contain an Euler path or circuit.
\end{CornellChecklist}

\end{document}
